\prefacesection{Lay Abstract}

Scientific software is supported by more than source code. Requirements,
design documents, equations, tests, and build instructions often repeat the
same facts in different forms. When people update each artifact separately,
the copies can disagree or important details can be omitted. This thesis
presents Drasil, a framework that stores scientific and engineering knowledge
once and uses it to generate consistent software and documentation. Drasil
represents concepts, symbols, units, equations, and relationships as reusable
``chunks,'' then uses ``recipes'' to select and arrange them for different
audiences. Reimplementing six scientific-computing case studies showed that
corrections made in one place propagated through generated artifacts, hidden
assumptions and inconsistencies became easier to detect, and closely related
software variants could reuse shared knowledge. The approach also has costs:
knowledge must be encoded carefully, and new contributors face a learning
curve. These results motivate better authoring tools and broader empirical
validation.
